(1) 含 $N$ 个原子的三维德拜固体，采用三支声学模近似且态密度 $g(\omega)=A\omega^2$（$0<\omega<\omega_D$）。由总模数 $3N$ 确定常数 $A$。
(2) 忽略零点能时建立声子的总内能积分并定义德拜温度 $\Theta_D$。
(3) 求 $T\ll\Theta_D$ 的热容首项及其 $T$ 的幂律。
(4) 求 $T\gg\Theta_D$ 时的杜隆–珀蒂热容，讨论低温与高温近似的适用范围。
